\documentclass[11pt, letterpaper]{article}

% Packages
\usepackage[utf8]{inputenc}
\usepackage{geometry}
\usepackage{graphicx}
\usepackage{listings}
\usepackage{xcolor}
\usepackage{amsmath}
\usepackage{float}
\usepackage{hyperref}
\usepackage{caption}

% Page Layout
\geometry{margin=1in}

% Code Listing Style
\definecolor{codegreen}{rgb}{0,0.6,0}
\definecolor{codegray}{rgb}{0.5,0.5,0.5}
\definecolor{codepurple}{rgb}{0.58,0,0.82}
\definecolor{backcolour}{rgb}{0.95,0.95,0.92}

\lstdefinestyle{verilogstyle}{
    backgroundcolor=\color{backcolour},   
    commentstyle=\color{codegreen},
    keywordstyle=\color{blue},
    numberstyle=\tiny\color{codegray},
    stringstyle=\color{codepurple},
    basicstyle=\ttfamily\footnotesize,
    breakatwhitespace=false,         
    breaklines=true,                 
    captionpos=b,                    
    keepspaces=true,                  
    numbers=left,                    
    numbersep=5pt,                  
    showspaces=false,                
    showstringspaces=false,
    showtabs=false,                   
    tabsize=2,
    language=Verilog
}

\lstset{style=verilogstyle}

% Title Info
\title{Homework 4: Pipelining Design}
\author{Yousef Awad \\ EEL 4783 - HDL in Digital Systems}
\date{\today}

\begin{document}

\maketitle

\section*{1)}
Translate the following C code into a Verilog code without pipelining.
\begin{lstlisting}[language=C]
int x = 0;
for (int i = 0; i < 3; i++)
{
  x = 2*x + 1;
}
\end{lstlisting}

\subsection*{Unpipelined Verilog Implementation}
\lstinputlisting{../src/unpipelined.sv}

\pagebreak
\section*{2) Analysis of Unpipelined Design}
Based on the implementation in \texttt{unpipelined.sv}, the following performance metrics are identified:

\begin{itemize}
    \item \textbf{Throughput:} 8 bits per clock cycle. Since the module produces one 8-bit result every clock cycle, the throughput is constant at 1 result/cycle.
    \item \textbf{Latency:} 1 clock cycle. The data enters the module and the final result of all three iterations is available at the next rising edge of the clock.
    \item \textbf{Timing (Critical Path Delay):} The critical path consists of three consecutive shift-and-add operations. 
    \[ T_{crit} = 3 \times T_{logic} + T_{co} + T_{su} \]
    Where $T_{logic}$ is the delay of a single $(2x+1)$ operation, $T_{co}$ is the clock-to-output delay of the input register, and $T_{su}$ is the setup time of the output register.
\end{itemize}

\section*{3) Pipelined Verilog Code}
\lstinputlisting{../src/pipelined.sv}

\pagebreak
\section*{4. Analysis of Pipelined Design}
Based on the implementation in \texttt{pipelined.sv}, the following performance metrics are identified:

\begin{itemize}
    \item \textbf{Throughput:} 8 bits per clock cycle. Once the pipeline is full (after the initial latency), the module produces one 8-bit result every clock cycle.
    \item \textbf{Latency:} 3 clock cycles. Since the design is divided into three stages, it takes three clock cycles for a specific input to propagate to the output.
    \item \textbf{Timing (Critical Path Delay):} The critical path is reduced to a single shift-and-add operation between registers.
    \[ T_{crit} = 1 \times T_{logic} + T_{co} + T_{su} \]
    This allows for a significantly higher maximum clock frequency compared to the unpipelined design, as the path delay is approximately one-third of the unpipelined version.
\end{itemize}

\end{document}
